\section{Taylor's theorem}
A Taylor polynomial matches a function's derivatives at its center. Matching
those derivatives does not by itself show that the infinite series represents
the function. The missing step is to show that the remainder tends to zero.
With Ross's convention the polynomial in $R_n$ stops at degree $n-1$:
\[
 R_n(x)=f(x)-\sum_{k=0}^{n-1}\frac{f^{(k)}(c)}{k!}(x-c)^k.
\]
Taylor's theorem bounds this error using the $n$th derivative at a point
between $c$ and $x$. Later exercises use Newton's method, which replaces a
curve by its tangent line and uses that line's zero as the next approximation.

\exercise{31.1}{The cosine series and its error}
\goal Find the Taylor series at $0$ and prove that it equals $\cos x$ everywhere.
\step{Compute the coefficients.}
Successive derivatives of cosine are $\cos x,-\sin x,-\cos x,\sin x$, then
repeat. At zero, the odd derivatives vanish and the $2k$th derivative is
$(-1)^k$. Thus the Taylor series is
\[
 \sum_{k=0}^{\infty}\frac{(-1)^k x^{2k}}{(2k)!}
 =1-\frac{x^2}{2!}+\frac{x^4}{4!}-\cdots.
\]
\step{Control the error, rather than assuming the series is correct.}
Every derivative of cosine has absolute value at most $1$. Taylor's theorem
therefore gives $\abs{R_n(x)}\le \abs{x}^n/n!$ for each fixed real $x$.
This bound tends to zero: if $a_n=\abs{x}^n/n!$, then
$a_{n+1}/a_n=\abs{x}/(n+1)\le1/2$ eventually (and $x=0$ is immediate).
The tail is consequently bounded by a geometric sequence tending to zero.
\answer{The displayed series equals $\cos x$ for every $x\in\R$.}

\exercise{31.2}{The hyperbolic sine and cosine series}
\goal Obtain series at zero for $\sinh x=(e^x-e^{-x})/2$ and
$\cosh x=(e^x+e^{-x})/2$, valid on all of $\R$.
\step{Add and subtract the exponential series.}
For each real $x$, the exponential series is absolutely convergent:
$e^x=\sum_{k\ge0}x^k/k!$ and $e^{-x}=\sum_{k\ge0}(-1)^kx^k/k!$.
Subtracting cancels the even terms; adding cancels the odd terms. Hence
\[
 \sinh x=\sum_{k=0}^{\infty}\frac{x^{2k+1}}{(2k+1)!},\qquad
 \cosh x=\sum_{k=0}^{\infty}\frac{x^{2k}}{(2k)!}.
\]
Termwise addition here is legitimate because it is addition of convergent
series at a fixed $x$.
\step{Check that these really are Taylor coefficients.}
Differentiation interchanges $\sinh$ and $\cosh$. Their values at zero are
$0$ and $1$, respectively. This gives exactly the displayed coefficients.
Alternatively, on $[-M,M]$ all derivatives of either function are bounded by
$e^M$, so Taylor's bound is $e^M\abs{x}^n/n!\to0$.
For any given $x$, choose $M>\abs{x}$; this proves validity at that $x$.
\answer{Both series represent their functions at every real number.}

\exercise{31.3}{Why the common derivative bound is unavailable}
\goal Explain the theorem choice for $f(x)=\log(1+x)$ in Example 2.
\step{Inspect the hypothesis of Corollary 31.4.}
That corollary requires a single constant bounding \emph{all} derivatives
on the interval in question. But
\[
 f^{(n)}(x)=(-1)^{n+1}(n-1)!(1+x)^{-n},\qquad
 \abs{f^{(n)}(0)}=(n-1)!.
\]
Even at the center $0$ these numbers have no common finite bound.
No interval containing zero can satisfy the corollary's hypothesis.
\step{Use the more precise remainder instead.}
For evaluating at $x=1$, Taylor's theorem gives a $y_n\in(0,1)$ with
\[
 \abs{R_n(1)}=\frac1{n(1+y_n)^n}\le\frac1n\longrightarrow0.
\]
The factorial growth of the derivative is canceled by the $n!$ in the
remainder. The corollary's stronger hypothesis would miss that cancellation.
\answer{Theorem 31.3 applies directly; the common-bound hypothesis of the
corollary fails.}

\exercise{31.4}{Smooth cutoffs and a smooth step}
\goal Construct the four requested infinitely differentiable functions.
\idea Start from the smooth function in Example 3,
\[
 \phi(t)=\begin{cases}e^{-1/t},&t>0,\\0,&t\le0.\end{cases}
\]
It is important that all its derivatives at zero vanish; this is what makes
the joining point smooth, despite the piecewise formula.
\step{(a) Move the start of the positive part to $a$.}
Set $f_a(x)=\phi(x-a)$. A translation preserves smoothness.
If $x\le a$ its argument is nonpositive, so $f_a(x)=0$; if $x>a$,
its argument is positive and $f_a(x)>0$.
\step{(b) Reverse the direction and move the endpoint to $b$.}
Set $g_b(x)=\phi(b-x)$. The affine change of variable preserves smoothness.
This is zero for $x\ge b$ and strictly positive for $x<b$.
\step{(c) Require both positivity conditions.}
Set $h_{a,b}(x)=f_a(x)g_b(x)$. Products of smooth functions are smooth.
Both factors are positive exactly when $a<x<b$. Outside this interval
at least one factor is zero, so the product is zero.
\step{(d) Turn the two cutoffs into a step.}
Set
\[
 h^*_{a,b}(x)=\frac{f_a(x)}{f_a(x)+g_b(x)}.
\]
The denominator never vanishes. Indeed, if $x\le a$ then $x<b$ and
$g_b(x)>0$; if $x\ge b$ then $x>a$ and $f_a(x)>0$; between $a$ and
$b$ both are positive. The permitted quotient rule for smooth functions
therefore applies. For $x\le a$ the quotient is $0$; for $x\ge b$
its numerator equals its denominator, so it is $1$.
\answer{These four explicit formulas have all the required properties,
including smoothness at $a$ and $b$.}

\exercise{31.5}{A nonzero smooth function with zero Taylor series}
\goal For $g(x)=e^{-1/x^2}$ when $x\ne0$, $g(0)=0$, find all derivatives
at zero and compare $g$ with its Taylor series.
\step{(a) Recognize a smooth composition.}
Using $\phi$ from Example 3, we have $g(x)=\phi(x^2)$ for every $x$.
Since both $\phi$ and $x\mapsto x^2$ are smooth, so is $g$.
Repeated chain and product rules show that each positive-order derivative
of this composition is a finite sum of polynomial factors times
$\phi^{(j)}(x^2)$, for integers $j\ge1$. At zero every such term vanishes,
because $\phi^{(j)}(0)=0$. Also $g(0)=0$. Thus
$g^{(n)}(0)=0$ for all $n=0,1,2,\ldots$.
\step{(b) Compare the function and the series point by point.}
Every coefficient $g^{(n)}(0)/n!$ is zero, so the Taylor series is the
identically zero series. It equals $g$ at $0$. For $x\ne0$,
$e^{-1/x^2}>0$, so it cannot equal $g(x)$ there.
\answer{Smoothness alone does not imply agreement with a Taylor series,
even arbitrarily close to its center.}

\exercise{31.6}{A telescoping proof of Taylor's theorem}
\goal Verify the older proof using its auxiliary function $F$.
Fix $x>0$. As in the theorem's proof, define $M$ by
\[
 f(x)=\sum_{k=0}^{n-1}\frac{x^k}{k!}f^{(k)}(0)+\frac{Mx^n}{n!}.
\]
\step{(a) Differentiate and cancel adjacent terms.}
Write the auxiliary function as
\[
 F(t)=\sum_{k=0}^{n-1}\frac{(x-t)^k}{k!}f^{(k)}(t)
       +M\frac{(x-t)^n}{n!},\qquad0\le t\le x.
\]
For $k\ge1$, the product rule gives
\[
 \frac{d}{dt}\left(\frac{(x-t)^k}{k!}f^{(k)}(t)\right)
 =-\frac{(x-t)^{k-1}}{(k-1)!}f^{(k)}(t)
   +\frac{(x-t)^k}{k!}f^{(k+1)}(t).
\]
The negative term cancels the positive term from index $k-1$.
Only the last positive term and the derivative of the $M$ term remain:
\[
 F'(t)=\frac{(x-t)^{n-1}}{(n-1)!}\bigl(f^{(n)}(t)-M\bigr).
\]
All needed derivatives exist on an open interval containing $[0,x]$;
in particular $F$ is continuous on the closed interval and differentiable
in its interior. For $n=1$ the formula follows directly from
$F(t)=f(t)+M(x-t)$.
\step{(b) Evaluate the endpoints.}
The definition of $M$ says $F(0)=f(x)$. At $t=x$ all positive powers
of $x-t$ vanish, leaving $F(x)=f(x)$. Thus $F(0)=F(x)$.
\step{(c) Apply Rolle's theorem.}
There is $y\in(0,x)$ with $F'(y)=0$. Since $x-y>0$, the prefactor
in the derivative formula is nonzero; hence $f^{(n)}(y)=M$.
Substituting into the defining equation for $M$ yields Taylor's remainder.
\answer{The auxiliary function proves the required mean-value expression
for the remainder.}

\exercise{31.7}{Newton's method for the square root of two}
\goal Recover the recursion from Exercise 9.5.
For $f(x)=x^2-2$, we have $f'(x)=2x$. Newton's update at a nonzero
$t_n$ is
\[
 t_{n+1}=t_n-\frac{t_n^2-2}{2t_n}
 =\frac{2t_n^2-t_n^2+2}{2t_n}
 =\frac{t_n^2+2}{2t_n}.
\]
With initial point $t_1=1$ this is exactly the earlier sequence.
Every iterate is positive, because a positive $t_n$ gives a positive
numerator and denominator. Thus no zero denominator is encountered.
\answer{The recursion is Newton's method for $x^2-2=0$, starting at $1$.}

\exercise{31.8}{Two real roots of a quartic}
\goal Count the real roots of $p(x)=x^4+x^3-1$ and approximate each to six decimals.
\step{(a) Establish the shape and count crossings.}
The derivative is $p'(x)=x^2(4x+3)$. It is negative for $x<-3/4$
and positive for $x>-3/4$, except for the isolated zero at $x=0$.
Thus $p$ is strictly decreasing on $(-\infty,-3/4]$ and strictly increasing
on $[-3/4,\infty)$. For strict increase across zero, apply the Mean Value
theorem separately on either side and combine the inequalities.
Now $p(-2)=7>0$, $p(-1)=-1<0$, $p(0)=-1<0$, and $p(1)=1>0$.
The Intermediate Value theorem supplies one root in $(-2,-1)$ and one
in $(0,1)$. Strict monotonicity on the two branches rules out any others.
\step{(b) Run the explicit update.}
Newton's formula is
\[
 x_{n+1}=x_n-\frac{x_n^4+x_n^3-1}{4x_n^3+3x_n^2}.
\]
Two runs, retaining guard digits, give
\[
\begin{array}{c|rr}
 n&x_0=-1.5&x_0=0.8\\\hline
 0&-1.500000000000&0.800000000000\\
 1&-1.398148148148&0.819758064516\\
 2&-1.380747193802&0.819173040661\\
 3&-1.380277903561&0.819172513397\\
 4&-1.380277569098&0.819172513396
\end{array}
\]
As a check independent of apparent convergence, direct substitution gives
opposite signs at the endpoints of each interval
$(-1.380277570,-1.380277568)$ and $(0.819172512,0.819172515)$.
Each interval contains the unique root on its branch and lies wholly
within the rounding range for the answer below.
\answer{The two roots, rounded to six decimal places, are
$-1.380278$ and $0.819173$.}

\exercise{31.9}{Solving the cosine fixed-point equation}
\goal Approximate the solution of $x=\cos x$ to six decimals.
Set $f(x)=x-\cos x$. Since $f'(x)=1+\sin x>0$ on $(0,\pi/2)$,
the root there is unique. Newton's iteration is
\[
 x_{n+1}=x_n-\frac{x_n-\cos x_n}{1+\sin x_n}.
\]
Starting at $x_0=1$ gives successively
\[
 0.750363867840,\quad0.739112890911,\quad
 0.739085133385,\quad0.739085133215.
\]
The signs of $f$ at $0.739085132$ and $0.739085135$ are negative
and positive, respectively. Continuity and strict increase bracket the
root, so the displayed digits are not justified solely by two similar iterates.
\answer{The root is $x=0.739085$ to six decimal places.}

\exercise{31.10}{Solving an exponential equation}
\goal Approximate the root of $xe^x=2$ in $(0,1)$ to six decimals.
Use $f(x)=xe^x-2$, with $f'(x)=(1+x)e^x>0$ on $(0,1)$.
Thus there is at most one root there; the opposite endpoint signs give
existence by continuity. Newton's update is
\[
 x_{n+1}=x_n-\frac{x_ne^{x_n}-2}{(1+x_n)e^{x_n}}.
\]
Starting with $x_0=1$ gives
\[
 0.867879441171,\quad0.852783373416,\quad
 0.852605526369,\quad0.852605502014.
\]
Direct substitution gives $f(0.852605501)<0<f(0.852605504)$,
so the unique root is in that small interval.
\answer{The six-decimal value is $0.852606$.}

\exercise{31.11}{A convergent Newton sequence must approach a zero}
\goal Under the stated bounded-derivative hypotheses, prove $f(x)=0$
when the Newton iterates converge to the interior point $x$.
\step{Rearrange the recursion before taking limits.}
Since $f'(x_{n-1})\ne0$, the update is defined and implies
\[
 f(x_{n-1})=(x_{n-1}-x_n)f'(x_{n-1}).
\]
Choose $B<\infty$ with $\abs{f'(t)}\le B$ throughout $(a,b)$.
Convergence of $(x_n)$ means both $x_n$ and $x_{n-1}$ tend to $x$,
so their difference tends to zero. Consequently
\[
 \abs{f(x_{n-1})}\le B\abs{x_{n-1}-x_n}\longrightarrow0.
\]
\step{Identify this limit using continuity.}
Differentiability at the interior point $x$ implies continuity there,
so also $f(x_{n-1})\to f(x)$. Uniqueness of a real limit gives $f(x)=0$.
We did not assume $f'$ was continuous or bounded away from zero;
the boundedness from above was the estimate actually used.
\answer{Every convergent sequence of these iterates with an interior limit
converges to a root.}

\exercise{31.12}{An interior point is no farther than an endpoint}
\goal If $a<c<b$, prove $\abs{c-x}\le\max\{\abs{a-x},\abs{b-x}\}$
for every real $x$.
\step{Case 1: $x\le c$.}
Then $c-x\ge0$ and $b-x>0$, so
$\abs{c-x}=c-x\le b-x=\abs{b-x}$.
\step{Case 2: $x>c$.}
Then $x-c>0$ and $x-a>0$, so
$\abs{c-x}=x-c\le x-a=\abs{a-x}$.
These cases exhaust the real line. In either case the relevant endpoint
distance is at most the maximum of the two endpoint distances.
\answer{The claimed inequality follows for all $x\in\R$.}
